\chapter{DASAR TEORI}

\section{Block Cipher dan SPN}
\textit{Block cipher} memetakan blok plainteks berukuran tetap menjadi
blok cipherteks dengan kunci rahasia. Untuk setiap kunci $K$, enkripsi
$E_K$ memiliki invers $D_K$ sehingga $D_K(E_K(P))=P$.
SPN mengulang lapisan substitusi, permutasi atau pencampuran, dan
penggabungan kunci. Substitusi nonlinier mendukung \textit{confusion},
sedangkan penyebaran pengaruh byte masukan mendukung \textit{diffusion}.

\Isi{Tambahkan rujukan teori confusion, diffusion, dan SPN yang benar-benar
digunakan. Pertahankan bab ini ringkas.}

\section{S-box, Permutasi, dan Operasi Tambahan}
S-box memetakan nilai byte menggunakan tabel. Pemetaan bijektif memiliki
tabel invers untuk dekripsi. Permutasi mengubah posisi byte tanpa mengubah
nilainya. Operasi XOR, rotasi bit, serta pencampuran pada $\mathrm{GF}(2^8)$
digunakan oleh implementasi Brainrot. Perkalian medan dalam kode memakai
polinomial reduksi $x^8+x^4+x^3+x+1$, yang juga digunakan
AES~\cite{aes}.

\section{Mode Operasi}
ECB, CBC, CFB, OFB, dan CTR adalah lima mode kerahasiaan yang dijelaskan
pada NIST SP 800-38A~\cite{mode}. ECB memproses blok secara independen;
CBC mengaitkan blok melalui cipherteks sebelumnya. CFB memakai umpan balik
cipherteks, OFB memakai umpan balik keluaran cipher, dan CTR mengenkripsi
nilai pencacah. Rincian implementasi dan penanganan blok akhir tersedia
pada Bab IV.

\section{Padding dan Verifikasi Data}
ECB dan CBC memerlukan masukan berupa blok lengkap. Untuk panjang blok
$b$ byte dan pesan $L$ byte, PKCS\#7 menambahkan $p=b-(L\bmod b)$ byte,
masing-masing bernilai $p$. Pesan yang sudah tepat kelipatan blok tetap
mendapat satu blok padding sehingga penghapusannya tidak ambigu~\cite{padding}.
Kesamaan byte hasil dekripsi dengan masukan membuktikan pemulihan data
pada pengujian tersebut. Deteksi perubahan cipherteks oleh pihak lain
memerlukan mekanisme verifikasi tambahan yang harus dijelaskan tersendiri.

\section{Ukuran Statistik}
Jika dua cipherteks $C$ dan $C'$ memiliki $N$ bit, persentase perubahan
bit dinyatakan sebagai
\begin{equation}
  A(C,C')=\frac{\operatorname{HW}(C\oplus C')}{N}\times100\%,
\end{equation}
dengan $\operatorname{HW}$ menyatakan banyaknya bit bernilai satu.
Untuk satu blok Brainrot, $N=128$. Entropi empiris byte dihitung dari
frekuensi relatif $p_v$ setiap nilai byte $v$:
\begin{equation}
  H=-\sum_{v=0}^{255}p_v\log_2 p_v,\qquad 0\log_2 0=0.
\end{equation}
Nilai maksimum distribusi byte adalah 8 bit/byte. Histogram menampilkan
jumlah kemunculan masing-masing nilai $v$. Ketiga ukuran ini memberikan
deskripsi statistik sampel; hasilnya perlu dibahas bersama ukuran sampel
dan perilaku mode operasi.
